% ============================================================
% Chapter 6: Conclusion and Future Work
% CT-MUSIQ Undergraduate Thesis
% ============================================================

\chapter{Conclusion and Future Work}
\label{chap:conclusion}

\section{Summary of Contributions}
\label{sec:conc_summary}

This thesis has presented CT-MUSIQ, a novel no-reference image quality assessment
framework for brain low-dose computed tomography, motivated by the pressing clinical
need for automated, scalable quality assurance in high-throughput LDCT pipelines.
The work spans the complete research cycle from literature review and problem
identification through architectural design, training methodology, systematic experimental
evaluation, and ablation validation.

The principal scientific and engineering contributions of this thesis are:

\textbf{1. CT-MUSIQ architecture.} A purpose-built NR-IQA system adapted from the
Multi-Scale Image Quality Transformer (MUSIQ) for the specific constraints of brain CT
imaging. CT-MUSIQ processes 193 patch tokens from a two-scale spatial pyramid
($224 \times 224$ and $384 \times 384$ pixels), encodes their spatial positions with a
hash-based factored embedding scheme, and processes the full token sequence with a
12-layer, 8-head Transformer encoder initialised from pretrained ViT-B/32 weights.

\textbf{2. Clinical preprocessing adaptation.} The brain soft-tissue windowing
transformation (level~40~HU, width~80~HU) is identified and motivated as a critical
preprocessing step that aligns the model's perceptual input with the radiologist's
viewing conditions during annotation---a methodological refinement not present in the
original LDCTIQAC challenge design, which was targeted at abdominal CT.

\textbf{3. Scale-consistency KL regularisation.} A novel composite loss function that
adds a Kullback-Leibler divergence term between per-scale quality prediction
distributions and the global quality prediction, enforcing cross-scale consistency and
providing effective regularisation on the small 700-image training set. Ablation studies
confirm this is the most impactful single contribution, providing +0.034 aggregate
performance improvement at the 30-epoch ablation budget.

\textbf{4. Two-stage fine-tuning with EMA and gradient clipping.} A carefully engineered
training protocol that stabilises the fine-tuning of a large pretrained Transformer on
a small medical dataset, including frozen encoder warm-up, linear learning rate warm-up,
cosine annealing, exponential moving average weight smoothing, and gradient clipping.

\textbf{5. Comprehensive benchmarking.} CT-MUSIQ is benchmarked against six published
LDCTIQAC 2023 leaderboard methods and the independently reimplemented Agaldran Combo
baseline, with systematic ablation studies confirming the individual contribution of each
proposed component.

\textbf{6. Open-source implementation.} A complete, reproducible PyTorch implementation
with centralised configuration management, deterministic seeding, and modular scripts
for training, evaluation, and ablation, enabling full reproduction of all reported
experiments.


\section{Conclusions}
\label{sec:conc_conclusions}

The experimental results lead to the following conclusions:

\textbf{Conclusion 1: Multi-scale Vision Transformers are an appropriate architecture
for LDCT NR-IQA.} CT-MUSIQ's Aggregate of 2.7235 places it among the top methods
evaluated on the LDCTIQAC 2023 benchmark, demonstrating that the multi-scale Transformer
paradigm---originally designed for natural image IQA---transfers effectively to the
medical CT domain with appropriate domain adaptations.

\textbf{Conclusion 2: Scale-consistency KL regularisation is the most impactful
proposed component.} Among all ablation configurations tested, the scale-consistency KL
loss at $\lambda = 0.10$ provides the largest single performance gain (+0.034 aggregate
at 30 epochs), validating the hypothesis that preventing per-scale overfitting is critical
for learning from 700 training samples.

\textbf{Conclusion 3: Ensemble diversity can outperform single-model regularisation on
small datasets.} The Agaldran Combo (Swin-T + ResNet-50 ensemble) achieves a higher
Aggregate (2.7682) than CT-MUSIQ (2.7235) despite having fewer parameters and simpler
training. This finding suggests that architectural diversity through ensembling may be
a more effective strategy than regularising a larger single model on datasets of this
scale.

\textbf{Conclusion 4: Classical NR-IQA methods are fundamentally inadequate for CT.}
The 0.60 aggregate gap between CT-MUSIQ and BRISQUE confirms the literature review
conclusion that Natural Scene Statistics-based methods fail on calibrated medical images
due to the domain mismatch between natural image statistics and Hounsfield Unit
distributions.

\textbf{Conclusion 5: Brain CT windowing alignment is a non-trivial methodological
choice.} The adaptation of preprocessing from the original challenge's abdominal
windowing to the clinically standard brain soft-tissue window is both clinically
motivated and technically significant, ensuring the model processes the same intensity
distribution as the annotating radiologists.


\section{Limitations and Self-Critique}
\label{sec:conc_limitations}

This thesis acknowledges the following significant limitations that should be considered
when interpreting the results:

\textbf{Limited dataset scope.} All experiments are conducted on a single dataset of
1,000 brain CT images from one challenge. Generalisation to different scanner models,
patient populations, dose levels, or anatomical regions is entirely unvalidated. A
clinically deployable system would require validation on multi-institutional, multi-vendor
datasets with diverse acquisition parameters.

\textbf{Small test set.} The 100-image test set provides insufficient statistical power
to claim statistically significant performance differences between methods with aggregate
scores differing by less than 0.1 points. The reported results are point estimates subject
to substantial sampling uncertainty.

\textbf{Non-identical baseline reimplementation.} The Agaldran Combo implemented here
is motivated by but not identical to the original challenge submission. The higher
performance of the reimplemented Agaldran Combo (2.7682 vs.\ published 2.7427) may
reflect the different training conditions rather than the architectural advantages.

\textbf{Hardware constraints.} The 6~GB VRAM limitation of the consumer GPU used in
this work restricts the experimental scope: only two scales can be used in full training,
batch size is limited to 10, and ensemble strategies for CT-MUSIQ itself (e.g., multiple
CT-MUSIQ models trained with different random seeds) could not be explored.

\textbf{No clinical validation.} CT-MUSIQ's quality scores are validated against
radiologist consensus labels from the LDCTIQAC 2023 challenge but have not been validated
in a real clinical setting---for instance, by demonstrating that high-quality scans
predicted by CT-MUSIQ have improved diagnostic yield for specific pathologies relative
to low-quality scans predicted by CT-MUSIQ.


\section{Future Work}
\label{sec:conc_future}

The findings of this thesis open several promising directions for future research:

\textbf{1. Multi-institutional dataset collection.} The most impactful direction for
clinical validation would be the collection and annotation of a large, multi-institutional
dataset covering diverse scanner types, patient demographics, dose levels, and anatomical
regions. Such a dataset would enable robust evaluation of generalisation and domain
adaptation strategies.

\textbf{2. CT-MUSIQ ensemble.} Combining multiple CT-MUSIQ models trained with different
random seeds or different KL regularisation weights would be expected to improve
aggregate performance through prediction diversity, similar to the benefit observed with
the Agaldran Combo. A five-model ensemble of CT-MUSIQ variants could potentially surpass
the Agaldran Combo baseline.

\textbf{3. CT-MUSIQ and Agaldran Combo ensemble.} Given the complementary strengths of
CT-MUSIQ (better PLCC calibration) and Agaldran Combo (better SROCC/KROCC), combining
their predictions through a learned or simple averaging strategy would be expected to
outperform either method alone.

\textbf{4. Three-scale CT-MUSIQ with extended training.} The ablation study showed that
the three-scale configuration (A6) slightly underperforms at 30 epochs but may benefit
from the full 100-epoch training. With access to a larger GPU, training the three-scale
variant to convergence would determine whether native-resolution ($512 \times 512$)
patches provide additional quality information beyond the two-scale pyramid.

\textbf{5. Vision-language model integration.} Emerging VLM-based approaches such as
IQAGPT \cite{chen2024iqagpt} demonstrate that combining visual quality features with
language-based reasoning can produce explainable quality assessments aligned with
radiologist diagnostic criteria. Integrating CT-MUSIQ's visual features with a
domain-adapted language model to generate textual quality reports alongside numerical
scores would significantly enhance clinical interpretability.

\textbf{6. Prospective clinical deployment study.} Embedding CT-MUSIQ within a hospital
PACS pipeline and conducting a prospective study measuring its agreement with
radiologist quality flags, its scan rejection accuracy, and its correlation with
diagnostic outcomes would provide the clinical validation necessary for regulatory
approval and widespread adoption.

\textbf{7. Explainability through attention visualisation.} CT-MUSIQ's Transformer
self-attention maps could be visualised to identify the image regions and patches that
most strongly influence the quality score prediction, providing a degree of explainability
that may help radiologists understand and trust the automated assessment.

\textbf{8. Extension to volumetric quality assessment.} The current work operates on
individual 2D axial slices. A volumetric extension that aggregates quality scores across
all slices of a CT volume---or processes slices jointly through a temporal/volumetric
attention mechanism---would better reflect clinical practice, where overall scan quality
is assessed at the volume level.

\textbf{9. Self-supervised pre-training on unlabelled CT.} Given the chronic shortage
of annotated CT quality data, self-supervised pre-training strategies adapted for CT
images (e.g., masked image modelling in HU space) could provide stronger CT-domain
representations as the initialisation for CT-MUSIQ, potentially surpassing the
ImageNet-pretrained ViT-B/32 starting point.


\section{Broader Impact and Ethical Considerations}
\label{sec:conc_ethics}

The deployment of automated AI systems in clinical medical imaging carries significant
ethical responsibilities that this section briefly examines, recognising that the thesis
work itself is an early research contribution rather than a clinically deployed system.

\subsection{Radiation Safety and the ALARA Principle}

CT-MUSIQ is motivated in part by the goal of supporting the ALARA principle: by providing
reliable automated quality assessment, the system could enable real-time scan rejection
and immediate re-acquisition at the dose level that achieves diagnostic quality, rather
than systematically exposing patients to higher doses ``to be safe.'' However, this
goal introduces a dual risk: if CT-MUSIQ incorrectly accepts a poor-quality scan
(false positive), a patient may receive a subdiagnostic examination that delays care;
if it incorrectly rejects a good-quality scan (false negative), the patient receives
unnecessary additional radiation. The appropriate operating threshold for a clinical
deployment must be determined through careful clinical validation studies that explicitly
quantify the trade-off between these two error types.

\subsection{Bias and Fairness}

The LDCTIQAC 2023 dataset, while a valuable benchmark, was collected under specific
acquisition conditions and annotated by a specific pool of radiologists. The model
trained on this data may inherit biases from the annotation process, including
systematic biases related to the specific scanner parameters used, the demographics
of the patient cohort, or the preferences of the annotating radiologists. Before
clinical deployment, it would be essential to evaluate CT-MUSIQ on data from diverse
patient populations, scanners, and annotation teams to assess the robustness and fairness
of its quality predictions across different subgroups.

\subsection{Transparency and Explainability}

A clinically deployed AI system must be explainable to the radiologists and
radiographers who interact with it. CT-MUSIQ, as a black-box neural network, does not
currently provide any explanation for its quality score predictions beyond the numerical
score itself. Attention visualisation techniques---such as attention rollout, Grad-CAM
applied to the \texttt{[CLS]} token, or attention map averaging across layers---could
provide heatmaps indicating which image regions most strongly influenced the quality
prediction. Such visualisations would help radiographers understand why a scan was
flagged as low quality, potentially identifying specific artefacts that the model found
problematic. Developing and validating these explainability tools is an important step
toward clinical trust and regulatory approval.

\subsection{Regulatory Pathway}

In most jurisdictions, AI-based clinical decision support tools (including automated
image quality assessment systems) require regulatory approval before clinical deployment.
In the European Union, such systems are classified as medical devices under the MDR
2017/745 and require conformity assessment by a notified body. In the United States,
FDA clearance is required under the 510(k) pathway for software as a medical device
(SaMD). These regulatory processes require extensive clinical validation evidence,
cybersecurity assessment, and quality management system certification. CT-MUSIQ, as a
research prototype, is far from meeting these requirements, but the work presented in
this thesis provides a scientifically rigorous foundation upon which a regulatory-grade
system could be built.

\section{Reflection on the Research Process}
\label{sec:conc_reflection}

This thesis was developed alongside the implementation of the CT-MUSIQ codebase,
creating a close integration between experimental work and documentation. Several
methodological lessons emerged that are worth recording for future researchers.

\textbf{Dataset inspection.} An early challenge was discovering that the dataset
contained brain CT slices rather than the abdominal CT implied by the original challenge
description. This finding---made through careful inspection of the raw images---led
to the critical methodological adaptation of using brain soft-tissue windowing, which
measurably improved model performance. Thorough dataset inspection is a prerequisite
for any medical AI project, and domain expertise is essential to recognising when
reported dataset metadata does not match the actual image content.

\textbf{Centralised configuration.} Maintaining a single \texttt{config.py} file as the
sole repository of all hyperparameters proved invaluable during the iterative training
and ablation process. When a hyperparameter needed adjustment, a single change propagated
correctly to all scripts, eliminating the risk of inconsistencies between training and
evaluation configurations. This practice is strongly recommended for all machine learning
research projects involving multiple interacting scripts.

\textbf{Training monitoring.} The training logs revealed unexpected patterns in the
validation aggregate that prompted investigation of the stochastic training dynamics.
Rather than being a cause for concern, these patterns were identified as normal behaviour
of gradient descent on a small dataset. Systematic monitoring of training dynamics
through detailed logging is essential for diagnosing model behaviour and distinguishing
normal stochasticity from genuine training pathologies.

\section{Closing Remarks}
\label{sec:conc_closing}

This thesis demonstrates that the challenge of automated perceptual quality assessment
for brain low-dose CT imaging can be meaningfully addressed through a carefully designed
combination of multi-scale Vision Transformer architecture, clinical domain adaptation,
and principled regularisation. CT-MUSIQ achieves competitive state-of-the-art performance
on the LDCTIQAC 2023 benchmark while operating entirely within the constraints of
consumer GPU hardware, making its principles accessible to researchers without access to
high-end computing infrastructure.

The work also reveals that ensemble methods combining architecturally diverse backbones
can surpass even well-regularised single models on small medical datasets, a finding with
broad implications for the development of deep learning systems in data-scarce clinical
domains. The complementary nature of CT-MUSIQ and the Agaldran Combo baseline suggests
that hybrid strategies combining multi-scale patch tokenisation with diverse
convolutional-hierarchical features represent a particularly promising direction for
future research.

Ultimately, the goal of this research is to contribute a step toward a future where every
CT acquisition is automatically and reliably assessed for diagnostic quality, enabling
real-time quality feedback to radiographers, objective benchmarking of reconstruction
algorithms, and---most importantly---improved patient care through the consistent delivery
of diagnostically acceptable images at the lowest achievable radiation dose.
